\documentclass[11pt]{article}
\usepackage[T1]{fontenc}
\usepackage{amsmath}
\usepackage{amssymb}
\usepackage{array}
\usepackage{booktabs}
\usepackage[margin=1in]{geometry}
\usepackage[hidelinks]{hyperref}
\usepackage{seqsplit}
\setlength{\emergencystretch}{2em}

\newcommand{\claimstatus}[2]{\textbf{#1 [#2]}}
\newcommand{\poly}{\operatorname{poly}}
\newcommand{\bits}{\operatorname{bitlength}}
\newcommand{\archivepaper}{\path{paper/main.tex}}
\newcolumntype{L}[1]{>{\raggedright\arraybackslash}p{#1}}
\newenvironment{claimproof}{\par\noindent\textit{Proof.}}{\hfill\(\square\)\par}


\title{Las Vegas Complete Factorization under Unrecognized\\
Hereditary Order Promises}
\author{MOSEF Research Program}
\date{31 July 2026}

\begin{document}
\maketitle

\begin{abstract}
Classical exponent--GCD factorization succeeds when prime components receive
different capped divisibility profiles, not merely when a schedule covers
their orders.  Using this separation semantics, we isolate two
complete-factorization theorems under hereditary order-asymmetry promises.
The first uses a fresh uniform base in each \(p-1\) schedule trial and splits
with probability at least \(5/12\) per complete witness cycle.  The second
uses fresh uniform Lucas parameters in a nonsplit \(p+1\) schedule and
succeeds with probability greater than \(1/12\) at a witness trial.  Exact
primality testing, maximal-perfect-power preprocessing, every modular
evaluation and GCD, verified recursion, and output are charged in the
standard bit length of the input.  The witness factors and promise
membership are neither supplied to nor recognized by the online algorithms.
With a finite per-node budget \(s\), total wrappers return either a complete
verified factorization or \texttt{UNRESOLVED} on every input.  On the two
promises their local unresolved tails are at most \((7/12)^s\) and
\((11/12)^s\), and their complete-recursion bounds are respectively
\(\min\{1,4m(7/12)^s\}\) and \(\min\{1,4m(11/12)^s\}\), without assuming
independence between recursive nodes.  Fresh sampling is essential: fixed
bases admit explicit hereditary counterexamples.  We further distinguish
divisor coverage from order-signature injectivity, show that conjugate Lucas
parameters do not create an independent channel, and record the finite
density limit of a common schedule.  No unproved number-theoretic conjecture
is assumed, but these are unrecognized-promise results, not a polynomial-time
algorithm for factoring arbitrary integers.
\end{abstract}

\begin{costmodel}{Cost model at a glance}
\costrow{Online charged}{Public schedule construction, exact preprocessing,
fresh sampling, modular \(p-1\) or Lucas evaluation, every GCD and divisor
check, recursion, and verified output.}
\costrow{Not supplied}{Witness prime factors, a promise-membership bit, and
factor-dependent schedules are not inputs and are not inferred.}
\costrow{Offline proof}{Group and root counts, tail bounds, counterexample
searches, and reproducibility records justify the theorem but are not online
advice.}
\costrow{Guarantee}{Expected polynomial bit complexity holds on the stated
hereditary promise; the bounded wrapper is total on all inputs and may return
\texttt{UNRESOLVED}.}
\end{costmodel}

\section{Scope and model}

For \(N\ge1\), write
\[
  m=\bits(N)=\lfloor\log_2N\rfloor+1.
\]
The input is binary, and every complexity statement applies the opening cost
box at each recursion node.  Exact primality testing and
maximal-perfect-power decomposition precede a splitter, and a reported
divisor is checked before its two children are processed.

Let \(B(k)\) and \(R(k)\) be fixed, nondecreasing, polynomially bounded
integer functions computable in \(\poly(k)\) time, and let
\[
  M_B(k)=\operatorname{lcm}(1,\ldots,B(k)).
\]
A composite non-perfect-power residual \(K\), of local length \(k\), has a
\(p-1\) witness when distinct \(p,q\mid K\) and \(1\le t\le R(k)\) obey
\[
  p-1\mid tM_B(k),\qquad q-1\nmid tM_B(k).
\]
The nonsplit \(p+1\) witness replaces these divisibilities by
\[
  p+1\mid tM_B(k),\qquad q+1\nmid tM_B(k),
\]
and is applied only to odd residuals.  A promise is \emph{hereditary} when
the corresponding witness exists at every composite non-perfect-power
residual reached by recursion.  These definitions depend on unknown prime
factors: the online algorithm is not given the witness and does not decide
membership.

\section{Fresh-base \texorpdfstring{\(p-1\)}{p-1} factorization}

\claimstatus{THM-001}{PROVED}.  For every fixed pair \(B,R\) above, fresh
uniform residues give a classical Las Vegas algorithm that completely
factors every input in the hereditary \(p-1\) semismooth-asymmetry promise,
is always correct, terminates almost surely, and has expected bit complexity
polynomial in the input length.  At a promised residual, one complete
schedule cycle splits with probability at least \(5/12\).

\begin{claimproof}
At residual \(K\), enumerate \(1\le t\le R(k)\).  For each \(t\), sample a
fresh exact uniform \(a\bmod K\), retain a proper \(\gcd(a,K)\), and otherwise
compute
\[
  \gcd\!\left(a^{tM_B(k)}-1,K\right).
\]
Fix a witness \(p,q,t\).  Every unit becomes one modulo \(p\).  The subgroup
of units modulo \(q\) killed by the exponent is proper, hence has at most
half the units.  Proper nonunits already expose a divisor, so success at the
witness trial has probability at least
\[
 \frac{K-1-\varphi(K)/2}{K}
 \ge \frac{K-1}{2K}\ge\frac5{12}.
\]
Fresh samples yield a geometric tail and almost-sure termination.  The
exponent has bit length
\[
O(B(k)\log B(k)+\log R(k)).
\]
Binary powering, GCD, exact
preprocessing, and the \(O(m)\)-node factor tree therefore give expected
\(\poly(m)\) bit cost.  Every output is a checked proper divisor, so the
algorithm is always correct.  The fully expanded branch and recursion proof
is archived at \path{research/proofs/THM-001-semismooth-promise.md}.
\end{claimproof}

Freshness cannot be dropped: for \(N=51\), base \(2\), and
\(d=\operatorname{lcm}(1,\ldots,8)=840\), the prime \(3\) is covered and
\(17\) is not, but \(\operatorname{ord}_{17}(2)=8\mid d\), producing the
full collision \(\gcd(2^{840}-1,51)=51\).

\section{Coverage is not separation}

For a finite positive exponent family \(\Delta\), define the divisibility
signature
\[
  \Sigma_\Delta(r)=\{d\in\Delta:r\mid d\}.
\]

\claimstatus{BAR-001}{PROVED}.  Divisor coverage does not imply order
separation.  A profile \((r_1,\ldots,r_s)\) is separated exactly when its
signatures are not all equal.  Every order in \(\{1,\ldots,n\}\) is separated
from every distinct order exactly when \(r\mapsto\Sigma_\Delta(r)\) is
injective, and an explicit nonempty injective signature family satisfies
\[
  |\Delta|\ge\lceil\log_2(n+1)\rceil.
\]

\begin{claimproof}
An exponent separates the profile precisely when it belongs to some but not
all of the sets \(\Sigma_\Delta(r_i)\); this is equivalent to unequal
signatures.  Injectivity is therefore necessary and sufficient for every
distinct pair.  Nonempty signatures occupy at most \(2^{|\Delta|}-1\)
subsets, proving the lower bound.  The smallest counterexample is
\(\Delta=\{2\}\): both orders one and two are covered but have the same
signature, and \(\gcd(5^2-1,6)=6\).
\end{claimproof}

\section{The conjugate channel is correlated}

Let \(V_0(P,1)=2\), \(V_1(P,1)=P\), and
\(V_{j+1}(P,1)=PV_j(P,1)-V_{j-1}(P,1)\).

\claimstatus{BAR-002}{PROVED}.  If \(a\) is a unit modulo \(N\) and
\(P=a+a^{-1}\), then for every \(d\ge0\),
\[
 V_d(P,1)-2=a^{-d}(a^d-1)^2,\qquad
 P^2-4=a^{-2}(a^2-1)^2\pmod N.
\]
Thus the derived Lucas residue has the same prime support as the
multiplicative residue.  When exponent two is retained, the derived channel
cannot enlarge the multiplicative family's success domain.

\begin{claimproof}
The sequence \(a^d+a^{-d}\) has the displayed initial values and recurrence,
so it equals \(V_d(P,1)\).  Subtraction of two and multiplication by the unit
\(a^d\) gives the square identity.  The discriminant identity is the case
\(d=2\).  Units and squares preserve prime support.  At prime powers,
valuations double and cap at the modulus valuation, so the map can worsen a
split; \((N,a,P,d)=(25,2,15,4)\) gives multiplicative GCD \(5\) but Lucas
GCD \(25\).
\end{claimproof}

\section{Independent nonsplit \texorpdfstring{\(p+1\)}{p+1} factorization}

\claimstatus{LEM-003}{PROVED}.  For every odd prime \(q\) and \(d>0\),
\[
 \#\{P\in\mathbb F_q:V_d(P,1)=2\}
 =\frac{\gcd(d,q-1)+\gcd(d,q+1)}2.
\]
Exactly \(1+\mathbf1_{2\mid d}\) roots are the degenerate parameters
\(P=\pm2\), and exactly \((q-1)/2\) parameters have nonsquare discriminant.

\begin{claimproof}
Parameterize a root by \(P=z+z^{-1}\).  Split \(z\) between
\(\mathbb F_q^\times\), of order \(q-1\), and the norm-one subgroup of
\(\mathbb F_{q^2}^\times\), of order \(q+1\).  The equality \(V_d(P,1)=2\)
is \(z^d=1\).  Each nondegenerate pair \(\{z,z^{-1}\}\) maps to one \(P\);
the fixed points \(z=\pm1\) contribute the stated correction.  Counting
roots in the two cyclic groups and simplifying yields the formula.  The
quadratic character count for \(P^2-4\) gives the nonsplit count.
\end{claimproof}

\claimstatus{THM-002}{PROVED}.  Fresh uniform parameters \(P\bmod K\) give a
classical Las Vegas complete-factorization algorithm on the hereditary
nonsplit \(p+1\)-asymmetry promise.  It is always correct, terminates almost
surely, and has expected polynomial bit complexity.  A witness trial
succeeds with probability greater than \(1/12\).

\begin{claimproof}
Retain a proper discriminant GCD \(\gcd(P^2-4,K)\), and independently retain
a proper sequence GCD
\(\gcd(V_d(P,1)-2,K)\), even when the discriminant GCD is full.  At the
witness prime \(p\), \((p-1)/2\) nonsplit parameters force the sequence to
two.  At \(q\), LEM-003 bounds the nondegenerate colliding parameters by
\[
 C_q(d)=\frac{\gcd(d,q-1)+\gcd(d,q+1)}2-2
 \le\frac{3q-9}{4}.
\]
CRT therefore gives probability
\[
 \frac{p-1}{2p}\frac{q-C_q(d)}q>\frac1{12}.
\]
Fresh parameters give the geometric tail.  Binary Lucas evaluation uses
the exponent's polynomial bit length, and the preprocessing and recursion
argument is the one used for THM-001.  The full proof is
\path{research/proofs/THM-002-nonsplit-lucas-promise.md}.
\end{claimproof}

\section{Bounded total wrappers}

Fix a positive integer \(s\). At each randomized composite node, replace the
unbounded cycle loop by at most \(s\) complete cycles. Prime,
maximal-perfect-power, even-factor, proper-factor, and exhausted-budget
branches are explicit. The only public outcomes are a complete prime list
whose product is the input and \texttt{UNRESOLVED}; partial factor lists are
discarded.

This bounded form is a total corollary of THM-001 and THM-002. It terminates
on every positive input and never returns a wrong factorization. On the
corresponding hereditary promise, conditional on reaching any randomized
node, the local unresolved probabilities obey
\[
 \Pr[U_{p-1}]\le\left(\frac7{12}\right)^s,\qquad
 \Pr[U_{p+1}]\le\left(\frac{11}{12}\right)^s.
\]
The complete recursion has fewer than \(4m\) invocations after maximal-power
preprocessing. A conditional union bound, which needs no independence
between recursive nodes, therefore gives
\[
 \Pr[U^{\mathrm{complete}}_{p-1}]
 \le\min\left\{1,4m\left(\frac7{12}\right)^s\right\},
\]
\[
 \Pr[U^{\mathrm{complete}}_{p+1}]
 \le\min\left\{1,4m\left(\frac{11}{12}\right)^s\right\}.
\]
Outside the promises only totality and no-wrong-factor correctness survive;
\texttt{UNRESOLVED} does not decide membership. The complete proof and exact
reference API are in
\path{research/proofs/M84-bounded-total-promise-wrappers.md} and
\path{python/mosef_reference/promise_wrappers.py}.

\section{A common-schedule limitation}

For an odd prime \(p\), define
\[
 \sigma_\Delta(p)=
 \bigl(\mathbf1_{p-1\mid d},\mathbf1_{p+1\mid d}\bigr)_{d\in\Delta}.
\]
This is a factor-dependent analytical object, not an input recognizer.

\claimstatus{BAR-003}{PROVED}.  For distinct odd primes \(p,q\), the union
of the local \(p-1\) and \(p+1\) asymmetry promises holds exactly when
\(\sigma_\Delta(p)\ne\sigma_\Delta(q)\).  On a finite prime set \(S\), if
\(s=|S|\), \(h\) signatures are nonzero, and \(D=\max\Delta\), the promised
pair fraction obeys
\[
 \rho_\Delta(S)\le
 \frac{h(2s-h-1)}{s(s-1)}
 \le \min\left\{1,\frac{8|\Delta|\sqrt D}{s-1}\right\}.
\]

\begin{claimproof}
Unequal signatures are exactly a coordinate that covers one prime side but
not the other.  At most \(h(2s-h-1)/2\) unordered pairs touch a nonzero
signature, giving the first bound.  Each \(d\) can be divisible by at most
\(\tau(d)\) values of \(p-1\) and at most \(\tau(d)\) values of \(p+1\);
\(\tau(d)\le2\sqrt d\) yields \(h\le4|\Delta|\sqrt D\) and the second bound.
The statement is finite-distribution and schedule-specific.
\end{claimproof}

\claimstatus{BAR-004}{PROVED}.  An explicit polynomial-size exponent list
whose maximum exponent bit length is \(o(k\log k)\) hits at most
\(2^{o(k)}\) odd primes through the \(p\pm1\) divisibility tests.  Hence its
covered-pair fraction vanishes on any separately stipulated
\(2^{\alpha k}\)-size common-length prime population.

\begin{claimproof}
Every hit has \(p-1\mid d\) or \(p+1\mid d\), so the number of hit primes is
at most twice the total divisor count.  The standard maximal-order bound
\(\log_2\tau(d)=O(\ell(d)/\log\ell(d))\), with
\(\ell(d)=o(k\log k)\), is \(o(k)\).  A polynomial number of exponents
therefore has \(2^{o(k)}\) total divisor support.  Substitution in the exact
pair bound of BAR-003 proves the conclusion.  This is not a lower bound for
all compact representations or adaptive schedules.
\end{claimproof}

\section{Primary-source positioning}

\paragraph{M83-R01: Pollard \(p-1\).}
Pollard's published procedure supplies the classical smooth-\(p-1\)
mechanism, including unit and full-collision branches
\cite{pollard1974theorems}. It does not state the hereditary promise, fresh
exact-uniform sampling at every trial, the \(5/12\) complete-cycle bound, or
the closed complete-recursion cost used in THM-001. THM-001 is therefore
presented as a scoped synthesis around an established mechanism, not as a
priority claim.

\paragraph{M83-R02: Williams \(p+1\).}
Williams supplies the Lucas-sequence mechanism and its split, nonsplit, and
zero-discriminant hypotheses \cite{williams1982pplusone}. LEM-003's exact
parameter-root count, the repository's treatment of every degeneracy branch,
and THM-002's greater-than-\(1/12\) witness bound are self-contained project
derivations. The audit assigns them no literature-priority status.

\paragraph{M83-R03: separating signatures.}
Pairwise distinction by binary membership signatures is classical
separating-system language \cite{katona1966separating}. BAR-001 and BAR-024
specialize that language to divisor/support coordinates, prime-power
valuations, and exact proper-GCD semantics. The specialization is the
paper's scope; binary signature injectivity itself is not claimed as a new
combinatorial concept.

The registered M83 matrix contains the full inspected-source comparison.
Failure to locate an identical formulation in this bounded audit is not
evidence of novelty or priority.

\section{Reproduction, limitations, and archival map}

The deterministic finite audits supporting the two probability proofs are
rerun by
\begin{verbatim}
python scripts/run_m3_semismooth_search.py
python scripts/check_m3_semismooth_differential.py
python scripts/run_m7_nonsplit_search.py
python scripts/check_m7_nonsplit_differential.py
python scripts/check_m84_promise_wrappers.py
python -m unittest discover -s tests -p test_promise_wrappers.py
\end{verbatim}
Their records are
\path{research/experiments/EXP-0003-m3-semismooth-search.md} and
\path{research/experiments/EXP-0006-m7-nonsplit-lucas.md}.  The experiments
check finite arithmetic and branches; they are not proofs of the probability
bounds.

The full formal proofs, counterexamples, cost ledgers, negative results, and
claim dependencies remain in \archivepaper,
\path{research/CLAIMS.md}, and \path{research/NEGATIVE_RESULTS.md}.
The machine-checked projection is
\path{schemas/m82-paper-portfolio-v1.json}.  Promise membership remains
unrecognized, no input-density theorem is supplied, and no claim here solves
general classical polynomial-time factorization.

\bibliographystyle{plain}
\small
\bibliography{paper/references}

\end{document}
